% Modernised reading — fonds Alexandre Grothendieck, shelfmark 29, batch 4
% (pages 61-80).
%
% This is an INTERPRETATION, not a transcription. It re-reads the pages in
% today's notation and vocabulary, states the hypotheses the manuscript leaves
% implicit, and drops the critical apparatus so that the mathematics reads
% straight through. Everything the brackets and underlines carried moves into a
% footnote. In any dispute of substance the transcription governs: whoever
% cites Grothendieck must cite batch-04.fr.tex.
%
% Scope: the folder was read WHOLE before any of this was written — all eleven
% batches, pages 1-216, in one pass — and one file was then written per batch.
%
% Departures from the page, each footnoted where it occurs:
%   - page 70 writes « Pic(X) -> Pic(U) » where its own reason (« comme Y est
%     factoriel ») and its own next diagram both require Pic(Y) -> Pic(U). The
%     true statement is written here;
%   - page 68 credits purity with the SURJECTIVITY of pi_1(U) -> pi_1(X), which
%     normality alone gives; purity is what makes it an ISOMORPHISM in
%     codimension >= 2, and that is where its own diagram puts it;
%   - the Frobenius argument of page 71 needs X and Y proper and connected over
%     an algebraically closed field for H^1(-,Z/p) to be the Frobenius
%     invariants of H^1(-,O); the hypothesis is supplied here;
%   - page 61 is a palimpsest of pencil under ink and is largely illegible. Its
%     theorem is not reconstructed: only what the transcription carries is used.
%
% Pass: Opus 5 (claude-opus-5), 2026-08-29 — first pass, unchecked against the
% pages by a human.
\documentclass[11pt,a4paper]{article}
% Le préambule commun à toutes les transcriptions du fonds Grothendieck.
%
% Il tient en une idée : **l'incertitude se compose, elle ne se résout pas.**
% Une transcription qui lisse un mot douteux a détruit la seule chose pour
% laquelle on la faisait — un lecteur doit pouvoir savoir, à chaque endroit,
% ce qui était sur le feuillet et ce qui a été ajouté. Les six macros qui
% suivent sont donc l'appareil critique, et non de la décoration.
%
% Les mêmes commandes sont reconnues par `scripts/render.mjs`, qui produit la
% vue de lecture du site. Une macro ajoutée ici doit l'être aussi là-bas,
% faute de quoi le rendu échouera bruyamment — ce qui est voulu : un rendu
% silencieusement faux est pire qu'un rendu absent.

\ProvidesPackage{grothendieck}[2026/08/08 Transcriptions du fonds Grothendieck]

\usepackage{amsmath,amssymb,amsthm}
\usepackage{mathrsfs}
\usepackage{stmaryrd}
% Les diagrammes commutatifs sont partout dans ce fonds : c'est la langue
% même de la géométrie algébrique de Grothendieck, et une transcription qui
% les rendrait en prose ne transcrirait rien.
\usepackage{tikz-cd}
% Français par défaut — c'est la langue des feuillets — mais l'anglais est
% chargé aussi : la traduction et le résumé sont des documents anglais, et la
% typographie française (espaces avant les deux-points) y serait une faute.
% Ces documents basculent par \selectlanguage{english} en tête de corps.
\usepackage[english,french]{babel}
\usepackage{fontspec}
\usepackage{geometry}
\geometry{a4paper,margin=25mm,marginparwidth=28mm}
\usepackage{marginnote}
\usepackage{xcolor}
% ulem, not soul. `soul`'s \st and \ul are built on a character-by-character
% scan that XeLaTeX's font machinery breaks: the document fails with an
% undefined control sequence rather than degrading. ulem does the same two jobs
% and is Unicode-engine safe. `normalem` stops it hijacking \emph, which the
% transcriptions use for Grothendieck's own emphasis.
\usepackage[normalem]{ulem}
\usepackage{hyperref}
\hypersetup{colorlinks=true,linkcolor=black,urlcolor=black,pdfborder={0 0 0}}

% --- Métadonnées du lot -----------------------------------------------------
% Reprises telles quelles par le script de rendu, qui les affiche en tête de
% la vue de lecture. Elles fixent ce que le fichier prétend couvrir : sans
% elles, une transcription flotte sans point d'attache dans un fonds de seize
% mille feuillets.

\newcommand{\folder}[1]{\gdef\grothFolder{#1}}
\newcommand{\batch}[1]{\gdef\grothBatch{#1}}
\newcommand{\leaves}[2]{\gdef\grothFirst{#1}\gdef\grothLast{#2}}
\newcommand{\foldertitle}[1]{\gdef\grothFoldertitle{#1}}
\gdef\grothFolder{?}\gdef\grothBatch{?}\gdef\grothFirst{?}\gdef\grothLast{?}\gdef\grothFoldertitle{}

% --- Le résumé --------------------------------------------------------------
%
% Il ouvre la lecture modernisée, sous le titre « Résumé », et il est composé
% en petit. Ce n'est pas un ornement : c'est le seul passage du document qui
% ne soit pas des mathématiques, et un lecteur doit pouvoir le reconnaître
% avant de l'avoir lu — puis le sauter s'il connaît déjà le sujet, ou s'y
% arrêter s'il ne le connaît pas. Le filet qui le suit marque la reprise du
% corps.

\newenvironment{resume}
  {\small\setlength{\parskip}{0.45em}}
  {\par\medskip\hrule\medskip}

% --- Mots-clés ---------------------------------------------------------------
%
% En anglais, en clôture du résumé de la lecture modernisée : le vocabulaire
% moderne sous lequel chercher ce que les feuillets construisent. Le manifeste
% du site (`npm run manifest`) les extrait pour étiqueter la cote entière —
% cette ligne est la source unique des tags, il n'y a pas de fichier à côté.
\newcommand{\keywords}[1]{\par\smallskip\noindent
  {\footnotesize\textsc{Keywords} --- \textit{#1}}\par}

% --- L'appareil critique ----------------------------------------------------

% Le numéro de feuillet, en marge. C'est l'ancre qui permet de régler un
% désaccord sur une formule en regardant *un* feuillet plutôt qu'une chemise
% de six cents. Le site s'en sert aussi pour tourner les pages du fac-similé
% quand on fait défiler la transcription.
\newcommand{\leaf}[1]{%
  \marginnote{\footnotesize\color{gray!70}\textsf{#1}}%
  \label{leaf:#1}\ignorespaces}

% Illisible : on ne devine pas. Un mot inventé qui se lit comme les autres est
% le pire résultat possible.
\newcommand{\ill}{\textcolor{red!60!black}{[\,\ldots\,]}}

% Lecture douteuse : le mot est donné, et signalé comme incertain.
\newcommand{\uncertain}[1]{\uline{#1}}

% Ajout de l'éditeur — une lettre manquante, un mot sous-entendu.
\newcommand{\add}[1]{\textcolor{blue!50!black}{[#1]}}

% Biffé par Grothendieck lui-même. Ce qu'il a rayé fait partie du document :
% une rature dit souvent où la pensée a changé de direction.
\newcommand{\struck}[1]{\sout{#1}}

% Note du transcripteur, sur l'état matériel du feuillet.
\newcommand{\note}[1]{\par\smallskip{\small\color{gray!60!black}\textit{[#1]}}\par\smallskip}

% Note marginale de Grothendieck, distincte du corps du texte.
\newcommand{\marginal}[1]{\par\smallskip\noindent\hspace{1em}%
  {\small\color{gray!60!black}#1}\par\smallskip}

% --- Titre ------------------------------------------------------------------

\AtBeginDocument{%
  \begin{center}
    {\large\bfseries Fonds Alexandre Grothendieck --- cote n\textsuperscript{o}~\grothFolder}\\[2pt]
    {\itshape\grothFoldertitle}\\[4pt]
    {\small feuillets \grothFirst--\grothLast\ (lot \grothBatch)}\\[2pt]
    {\footnotesize\color{gray!60!black}Transcription établie d'après le fac-similé numérisé
     par l'Université de Montpellier.\\
     Les lectures incertaines sont soulignées, les illisibles notées [\,\ldots\,],
     les ajouts entre crochets.}
  \end{center}
  \vspace{1em}}

\endinput


\folder{29}
\batch{4}
\pages{61}{80}
\dating{1959-1969}
\watermark{Édition de démonstration\\interprétation personnelle de l'œuvre}
\foldertitle{Groupe fondamental [Autour de SGA 4 et SGA 7] — lecture modernisée}

\begin{document}

\section*{Résumé}

\begin{resume}
Le lot précédent posait un programme de finitude et en démontrait les lemmes.
Celui-ci en tire les théorèmes, et il le fait deux fois : une fois pour le
groupe fondamental, une fois pour le groupe de Picard. Les deux démonstrations
se ressemblent au point qu'on peut les lire d'une seule traite, et
Grothendieck les a paginées à la suite, A à E puis F à L, dans une même
chemise.

Le procédé est le suivant, et il est classique de forme. Pour majorer un
invariant d'une variété de dimension quelconque, on la coupe par un espace
linéaire \emph{générique} jusqu'à tomber sur une courbe, et l'on montre que
l'invariant de la courbe majore celui de la variété. Ce qui rend l'opération
licite est le lemme de Bertini de Zariski : une section linéaire générique
d'un schéma irréductible reste irréductible. Le corollaire qu'en tire
Grothendieck est le pivot du lot : si $X$ est normal, l'application du groupe
fondamental de la section générique dans celui de $X$ est \emph{surjective}.
La courbe obtenue est projective normale, de genre $g$ ; son groupe fondamental
est engendré par $2g$ éléments ; donc celui de $X$ aussi.

Ensuite vient un pas que le carton n'avait pas encore fait, et qui est le
plus intéressant. Une extension de corps de type fini $K/k$ n'a pas de modèle
privilégié : elle en a une infinité, et ils se dominent les uns les autres.
Grothendieck passe à la limite sur tous les modèles propres et normaux, et
appelle \emph{groupe fondamental birationnel} de $K/k$ ce que la limite donne.
Le théorème de pureté fait que la limite est raisonnable, et le théorème dit
que ce groupe est encore topologiquement de type fini. Le même passage à la
limite est fait pour le groupe de Picard, avec un résultat plus fort encore :
le système est \emph{essentiellement constant}, c'est-à-dire que le groupe de
Picard ne dépend pas du modèle --- ce qu'on appelle l'invariance birationnelle
de $\mathrm{Pic}^{\tau}$.

Deux détails valent d'être signalés à qui lit ces pages pour la première fois.
Le premier est le seul endroit du carton où Grothendieck écrive un argument de
caractéristique $p$ pour lui-même : il découpe $H^{1}(X, \mathcal{O})$ en deux
morceaux selon le Frobenius, montre que l'un est de la cohomologie étale à
coefficients $\mathbb{Z}/p$ --- donc affaire de groupe fondamental --- et que
l'autre est affaire de normalité. Le second est la question qu'il pose en
dernier, et qui reste ouverte sur la page : cette majoration
a-t-elle une démonstration \emph{directe}, qui ne passe ni par le groupe
fondamental ni par la propreté du schéma de Picard ? En marge, trois points
d'interrogation.

Le lot se termine sur le retour de Murre, deux ans après : une lettre du 30
avril 1969 accompagnant la première version de ses notes sur les revêtements
formels étales et le théorème de Mumford, avec deux pages de commentaires
chapitre par chapitre. Le lot 5 contient les notes de lecture que
Grothendieck en a prises.

\keywords{birational fundamental group, Bertini theorem, Picard scheme, Frobenius, purity theorem}
\end{resume}

\pagerange{61}{64}
\section*{Fin du théorème de finitude relatif : deux corollaires sur les produits (pages 61 à 64)}

Les pages 61 à 63 portent, dans l'angle supérieur gauche, les numéros 2, 3 et 4
que Grothendieck donne à la suite ouverte page 46. La page 61 est un
palimpseste --- un premier jet au crayon repris à l'encre, puis biffé --- et sa
plus grande partie est perdue ; cette lecture n'y supplée pas. Ce qui en reste
lisible est la réduction à une courbe : $X'$ le normalisé de $X_{\mathrm{red}}$,
$D_{i}$ les composantes irréductibles de codimension 1 de l'image inverse d'un
diviseur $D$, $U = X - D$, et le passage du triplet $(k, X, U)$ au triplet
$(K, C, V)$ où $C$ est une \emph{courbe}.

Deux corollaires en sortent, et ils sont propres.

\textbf{Corollaire (produits de paires).} Soient $X'$, $X''$ propres sur $k$,
et $Z' \subset X'$, $Z'' \subset X''$ des parties fermées. Posons
\[
  X = X' \times X'' , \qquad
  Z = (Z' \times X'') \cup (X' \times Z'') .
\]
Alors, pour des points-base compatibles,
\[
  \pi^{t}_{1}(X ; Z) \;\xrightarrow{\ \sim\ }\;
  \pi^{t}_{1}(X' ; Z') \times \pi^{t}_{1}(X'' ; Z'')
\]
est un \emph{isomorphisme}. Autrement dit : le groupe fondamental modéré d'un
produit de paires est le produit des groupes fondamentaux modérés.

\textbf{Corollaire (produits, partie première à $p$).} Soient $X$, $Y$ de type
fini sur $k$. Alors
\[
  \pi'_{1}(X \times Y) \;\xrightarrow{\ \sim\ }\;
  \pi'_{1}(X) \times \pi'_{1}(Y) .
\]

Cette dernière est l'énoncé que la lettre à Serre de 1959 conjecturait, à la
page 128 du même carton, comme conséquence de sa suite exacte d'homotopie ---
« elle impliquerait que pour des schémas complets $X$, $Y$ au-dessus d'un
corps algébriquement clos $k$, on a
$\pi_{1}(X\times Y/k) = \pi_{1}(X/k)\times\pi_{1}(Y/k)$ ». Dix ans, et une
centaine de pages du carton, séparent la conjecture du
corollaire.\footnote{À une différence près, qui est celle de tout le carton :
la lettre de 1959 énonce le produit pour le groupe fondamental \emph{total},
partie infinitésimale comprise --- un pro-schéma en groupes --- et la page 62
l'énonce pour le quotient premier à la caractéristique. Ce n'est pas le même
énoncé, et c'est la restriction que le lot 7 rend intelligible.}

La restriction à $\pi'_{1}$ est signalée sur la page même : « mais elles \dots\
la car. ». C'est l'aveu que le $p$-groupe fondamental d'un produit n'est pas le
produit des $p$-groupes fondamentaux --- ce qui est exact, et c'est l'un des
phénomènes propres à la caractéristique $p$.

La page 63 énonce une variante de la question. Soit $X'$ un normalisé de $X$,
$Y \subset X$ fermé, $U = X - Y$ ; soient $Y_{1}, \dots, Y_{r}$ les composantes
de $Y$ de codimension 1 dans $X$ et $Y'_{ij}$ leurs images inverses dans $X'$.
Considérons les revêtements étales de $X - Y$ qui deviennent étales au-dessus
des anneaux locaux des $Y'_{ij}$ : ils sont classifiés par un groupe de type
galoisien $\pi$.

\textbf{Théorème.} Le groupe $\pi$ est topologiquement de génération finie.

Deux corollaires : pour $X$ complète et $Y$ de codimension $\geqslant 2$, un
quotient convenable de $\pi_{1}(X - Y)$ est topologiquement de génération
finie ; et pour $X$ algébrique, $\pi'_{1}(X)$ est topologiquement de type
fini.\footnote{Le corollaire 1 porte, à l'endroit du quotient, un mot que la
transcription laisse illisible. Cette lecture ne le devine pas : le groupe
quotient n'est pas nommé, et l'énoncé reste incomplet tel que la page le
laisse.}

La page 64 est un brouillon : des carrés de schémas, la surjectivité de
$\pi_{1}(X\times Y)\to\pi_{1}(X)\times\pi_{1}(Y)$, et en marge « surj ?? » en
regard de $\pi_{1}(Y)\to\pi_{1}(X)$ pour $X$ normal. Elle n'est pas reprise
ici : rien n'y est énoncé.

\pagerange{65}{69}
\section*{« $\pi_{1}$ birationnel » (pages 65 à 69)}

Grothendieck pagine cette suite A à E, et la suivante F à L.

\subsection*{Le théorème}

\textbf{Théorème.} Soient $k$ algébriquement clos et $X$ un schéma connexe sur
$k$ dont les composantes irréductibles $X_{i}$ soient isomorphes --- avec leur
structure réduite induite --- au complémentaire d'une partie fermée de
codimension $\geqslant 2$ dans un schéma propre sur $k$. Alors $\pi_{1}(X)$ est
topologiquement de type fini.

\emph{Démonstration.} Par SGA IX 6.2 on se ramène au cas où $X$ est intègre,
donc de la forme $X = Z - Y$ avec $Z$ propre et $Y$ de codimension $\geqslant 2$
dans $Z$ ; on peut supposer $Z$ intègre. Si $Z' \to Z$ est propre surjectif avec
$Z'$ intègre, et $X'$ l'image inverse de $X$, il suffit de prouver l'énoncé pour
$X'$ ; par le lemme de Chow on peut donc supposer $Z$ \emph{projectif}.

Récurrence sur $\dim Z$. Si $\dim Z \leqslant 0$, trivial. Si $\dim Z = 1$,
alors $Y = \varnothing$ puisque $Y$ est de codimension $\geqslant 2$, donc
$X = Z$ et le théorème est connu. Si $\dim Z \geqslant 2$, on utilise le lemme
suivant.

\subsection*{Le lemme de Bertini de Zariski}

\textbf{Lemme.} Soient $k$ un corps, $f : X \to \mathbf{P}^{r}_{k}$ un
morphisme avec $X$ \emph{irréductible} et $\dim \overline{f(X)} \geqslant n$.
Soit $1 \leqslant s \leqslant n-1$, soit $u$ le point générique d'un produit de
$s$ facteurs paramétrant les hyperplans, $K$ une clôture algébrique de $k(u)$,
$L_{u}$ la variété linéaire de $\mathbf{P}^{r}_{K}$ définie par $u$, et
\[
  X_{u} \;=\; X \times_{\mathbf{P}^{r}_{k}} L_{u}
        \;=\; X_{K} \times_{\mathbf{P}^{r}_{K}} L_{u} .
\]
Alors $X_{u}$ est irréductible.

En marge, la question et sa réponse, de sa main : « peut-on remplacer
\emph{irréductible} par \emph{connexe} ? --- non ». C'est un avertissement, non
une remarque : c'est l'irréductibilité, non la connexité, qui se transmet à une
section générique, et l'argument suivant repose entièrement sur cette
différence.

\textbf{Corollaire.} Si $X$ est \emph{normal}, alors
$\pi_{1}(X_{u}) \to \pi_{1}(X)$ est \emph{surjectif}.

\emph{Démonstration.} Soit $X' \to X$ un revêtement étale connexe. Puisque $X$
est normal, $X'$ est normal, donc connexe implique irréductible ; par le lemme,
$X'_{u} = X' \times_{X} X_{u}$ est irréductible, donc connexe. Un revêtement
étale connexe de $X$ reste donc connexe au-dessus de la section générique, ce
qui est exactement la surjectivité annoncée.

\begin{tikzcd}[column sep=small, row sep=small, nodes={font=\scriptsize}]
X'_u \arrow[r] \arrow[d] & X' \arrow[d] \\
X_u \arrow[r] \arrow[d] & X_K \arrow[r] \arrow[d] & X \arrow[d] \\
L_u \arrow[r] & \mathbf{P}^r_K \arrow[r] & \mathbf{P}^r_k
\end{tikzcd}

\subsection*{La conclusion : $2g$ générateurs}

Reprenons $Z$ normal de dimension $n \geqslant 2$ et $X = Z - Y$. Plongeons $Z$
dans $\mathbf{P}^{r}$, et prenons $s = n-1$, de sorte que $L_{u}$ soit de
codimension $n-1$. Alors
\[
  X_{u} \;=\; X_{K} \cap L_{u} \;=\; Z_{K} \cap L_{u} \;\neq\; \varnothing ,
\]
la seconde égalité parce que $Y_{K} \cap L_{u} = \varnothing$ --- $Y$ étant de
dimension $\leqslant n-2$ et $L_{u}$ générique de codimension $n-1$. Donc
$X_{u}$ est une courbe projective, d'ailleurs normale. Si $g$ en est le genre,
le corollaire au lemme de Bertini donne : \emph{$\pi_{1}(X)$ admet $2g$
générateurs topologiques}.\footnote{Le groupe fondamental d'une courbe
projective lisse de genre $g$ sur un corps algébriquement clos est
topologiquement engendré par $2g$ éléments, en toute caractéristique ---
en caractéristique $p$ c'est un théorème, obtenu par spécialisation depuis la
caractéristique zéro, et il n'y a pas de présentation finie du groupe entier :
les relations manquent. C'est exactement la différence entre l'énoncé a) et
l'énoncé b) de la page 40 (« de type fini » contre « de présentation finie »).}

\subsection*{Le passage à la limite sur les modèles}

Soit $K$ une extension de type fini de $k$. Les modèles propres et normaux de
$K$, ordonnés par domination, forment un ensemble ordonné filtrant croissant.

Soit $X$ un tel modèle et $U \subset X$ un ouvert. Deux faits, à distinguer
soigneusement :

\begin{enumerate}
  \item si $X$ est \emph{normal} et $U$ est un ouvert dense, l'homomorphisme
    $\pi_{1}(U) \to \pi_{1}(X)$ est \emph{surjectif} --- un revêtement étale
    connexe de $X$ reste connexe au-dessus de $U$ ;
  \item si $X$ est \emph{régulier} et $X - U$ de codimension $\geqslant 2$,
    l'homomorphisme est un \emph{isomorphisme} : c'est le \emph{théorème de
    pureté}.
\end{enumerate}

C'est cette distinction que le diagramme de la page 68 porte, l'isomorphisme du
bas étant annoté « th.\ de pureté ».\footnote{La phrase de la page, elle,
attribue la surjectivité à la pureté : « L'injection induite $U \to X$ étant le
complémentaire d'une partie fermée de codim.\ $\geqslant 1$, on voit que (par
pureté) on a un homo.\ \emph{surjectif} (si $X$ normal) ». En codimension 1 la
pureté ne dit rien, et la normalité suffit ; en codimension $\geqslant 2$ la
pureté donne l'isomorphisme, et c'est bien là que le diagramme la place. La
présente lecture sépare les deux énoncés.}

\begin{tikzcd}
\pi_1(X_1) \arrow[d, "f_{*}"'] & \pi_1(U_1) \arrow[l] & \\
\pi_1(X_0) & \pi_1(U_0) \arrow[l] \arrow[r, "\sim"] \arrow[ul, "g_{*}"] & \pi_1(U_1)
\end{tikzcd}

À la limite, on obtient un homomorphisme \emph{surjectif}
\[
  \pi_{1}(U, \Omega) \longrightarrow
  \varprojlim_{X \text{ modèle propre}} \pi_{1}(X, \Omega) ,
\]
$\Omega$ étant une extension algébriquement close de $K$. Le terme de droite est
ce que Grothendieck appelle le \emph{groupe fondamental birationnel} de $K/k$,
et qu'il note $\pi_{1}(K/k, \Omega)$.

\textbf{Théorème.} Si $K/k$ est une extension de type fini avec $k$
algébriquement clos, alors $\pi_{1}(K/k, \Omega)$ est topologiquement de type
fini.\footnote{Le mot \emph{birationnel} est de lui : « Le dernier est ce que
j'ai appelé le \emph{groupe fondamental birationnel} de $K/k$ ». Le passage à
la limite sur les modèles propres normaux d'un corps de fonctions est
aujourd'hui un objet standard, et la finitude est ce que la borne $2g$ ci-dessus
donne : le genre de la courbe générique ne dépend pas du modèle choisi, et
c'est le même entier qui majorera le rang de $\mathrm{Pic}$ à la page 76.}

\pagerange{70}{76}
\section*{« Pic birationnel » (pages 70 à 76)}

\subsection*{Le théorème}

\textbf{Théorème.} Soit $f : X \to Y$ un morphisme dominant \emph{birationnel}
de préschémas intègres normaux, propres sur $k$, avec $Y$ localement régulier.
Alors
\[
  f^{*} : \underline{\mathrm{Pic}}^{\tau}_{Y} \longrightarrow
  \underline{\mathrm{Pic}}^{\tau}_{X}
\]
est un \emph{isomorphisme}.

\emph{C'est un monomorphisme} parce que $f_{*}(\mathcal{O}_{X}) =
\mathcal{O}_{Y}$ universellement --- ce qui est le théorème principal de
Zariski, $f$ étant propre birationnel et $Y$ normal --- donc une immersion
fermée, s'agissant d'un homomorphisme de schémas en groupes propres.

\emph{Montrons qu'il est surjectif.} Soit $U$ l'ouvert de $Y$ où $f^{-1}$ est
défini ; son complémentaire est de codimension $\geqslant 2$. Comme $Y$ est
localement régulier, donc factoriel, la restriction
\[
  i^{*} : \mathrm{Pic}(Y) \xrightarrow{\ \sim\ } \mathrm{Pic}(U)
\]
est un isomorphisme.\footnote{La page écrit ici
« $\mathrm{Pic}(X) \xrightarrow{\ \sim\ } \mathrm{Pic}(U)$ », alors que sa
propre raison --- « comme $Y$ est factoriel » --- et son propre diagramme,
qui porte $i^{*} : \mathrm{Pic}(Y) \to \mathrm{Pic}(U)$, demandent l'un et
l'autre $\mathrm{Pic}(Y)$. C'est un lapsus, et l'énoncé écrit ici est celui que
l'argument utilise ; sans lui la factorialité de $Y$ ne servirait à rien.}
Soit $g : U \to X$ la section de $f$ sur $U$ et $i : U \to Y$ l'inclusion, de
sorte que $i = f \circ g$ et $i^{*} = g^{*} \circ f^{*}$ :

\begin{tikzcd}
& X \arrow[d, "f"] \\
U \arrow[ur, "g"] \arrow[r, "i"'] & Y
\end{tikzcd}

\begin{tikzcd}
& \mathrm{Pic}(X) \arrow[dl, "g^{*}"] \\
\mathrm{Pic}(U) & \mathrm{Pic}(Y) \arrow[u, "f^{*}"'] \arrow[l, "i^{*}"]
\end{tikzcd}

Puisque $i^{*}$ est un isomorphisme, $f^{*}$ est injectif et admet
$(i^{*})^{-1} g^{*}$ pour rétraction : $\mathrm{Pic}(Y)$ est un
\emph{facteur direct} de $\mathrm{Pic}(X)$, d'image l'ensemble des $\xi$ tels
que $\xi = f^{*}(i^{*})^{-1}g^{*}(\xi)$.

Le supplémentaire est engendré par les diviseurs $\sum n_{i} X_{i}$ portés par
les composantes exceptionnelles de $f$. Il reste donc à voir qu'une telle
combinaison $\tau$-équivalente à zéro est nulle --- « c'est là un lemme bien
connu », et c'est en effet la non-dégénérescence de la forme d'intersection sur
les diviseurs exceptionnels. Ainsi $f^{*}$ est une immersion \emph{surjective},
et il ne reste plus qu'à voir qu'elle induit un isomorphisme sur les espaces
tangents à l'origine, c'est-à-dire le

\textbf{Corollaire.} Sous les conditions du théorème,
\[
  f^{*} : H^{1}(Y, \mathcal{O}_{Y}) \longrightarrow H^{1}(X, \mathcal{O}_{X})
\]
est un isomorphisme.

\subsection*{Le corollaire en caractéristique $p$}

L'injectivité est acquise ; c'est la surjectivité qu'il faut. En caractéristique
$p > 0$, il suffit de prouver que sont bijectifs les deux homomorphismes
\[
  H^{1}(Y, \mathcal{O}_{Y})^{F} \longrightarrow H^{1}(X, \mathcal{O}_{X})^{F} ,
  \qquad
  {}_{F}H^{1}(Y, \mathcal{O}_{Y}) \longrightarrow
  {}_{F}H^{1}(X, \mathcal{O}_{X}) ,
\]
où $F$ est le Frobenius : les invariants d'une part, la partie tuée par $F$
d'autre part.

Le premier s'écrit, par la suite exacte d'Artin--Schreier
$0 \to \mathbb{Z}/p \to \mathcal{O} \xrightarrow{F-1} \mathcal{O} \to 0$,
\[
  H^{1}(Y, \mathbb{Z}/p\mathbb{Z}) \longrightarrow
  H^{1}(X, \mathbb{Z}/p\mathbb{Z}) ,
\]
et c'est un isomorphisme parce que $\pi_{1}(X) \to \pi_{1}(Y)$ en est un, par
le théorème de pureté sur $Y$ régulier.\footnote{L'identification demande
$X$ et $Y$ propres et connexes sur un corps algébriquement clos : la suite
d'Artin--Schreier donne
$0 \to H^{0}(\mathcal{O})/(F-1) \to H^{1}(\mathbb{Z}/p) \to
H^{1}(\mathcal{O})^{F} \to 0$, et le terme de gauche s'annule puisque
$H^{0}(\mathcal{O}) = k$ et que $F-1$ est surjectif sur un corps
algébriquement clos. La page écrit l'identification sans l'hypothèse ; elle est
supplée ici.}

Le second s'écrit, $X$ et $Y$ étant réduits,
\[
  H^{0}\bigl(Y, \mathcal{O}_{Y}/\mathcal{O}_{Y}^{p^{\nu}}\bigr) \longrightarrow
  H^{0}\bigl(X, \mathcal{O}_{X}/\mathcal{O}_{X}^{p^{\nu}}\bigr) ,
\]
ou encore, $k$ étant algébriquement clos donc parfait,
\[
  f^{*} : H^{0}\bigl(Y, \mathcal{O}_{Y}^{p^{-\nu}}/\mathcal{O}_{Y}\bigr)
  \longrightarrow
  H^{0}\bigl(X, \mathcal{O}_{X}^{p^{-\nu}}/\mathcal{O}_{X}\bigr) .
\]
Il est bijectif grâce à la \emph{normalité} de $X$ et de $Y$ : avec les
notations ci-dessus, on a le même triangle, où $g^{*}$ est injectif parce que
$\mathcal{O}_{X}^{p^{-\nu}}/\mathcal{O}_{X}$ est sans torsion ; donc $f^{*}$ et
$(i^{*})^{-1}g^{*}$ sont réciproques l'un de l'autre.

\begin{tikzcd}
& T(X) \arrow[dl, "g^{*}"] \\
T(U) & T(Y) \arrow[u, "f^{*}"'] \arrow[l, "i^{*}"]
\end{tikzcd}

C'est le seul endroit du carton où le découpage par le Frobenius soit écrit
pour lui-même, et il vaut d'être retenu comme méthode : \emph{la partie de
$H^{1}(\mathcal{O})$ fixée par le Frobenius est du groupe fondamental, la
partie tuée par lui est de la normalité}.

\subsection*{Le corollaire en caractéristique zéro}

On peut supposer que $k$ est le corps des fonctions d'un schéma intègre $S$ de
type fini sur $\mathbb{Z}$, et que $f : X \to Y$ provient d'un morphisme de
$S$-préschémas $\bar{f} : \bar{X} \to \bar{Y}$, intègres, $\bar{Y}$ lisse sur
$S$, $\bar{X}$ normal sur $S$, à fibres connexes. On peut alors supposer $X$ et
$Y$ plats sur $S$, et l'on a l'injectivité de
$R^{1}q_{*}(\mathcal{O}_{\bar{Y}}) \to R^{1}p_{*}(\mathcal{O}_{\bar{X}})$ ; la
constance de la dimension sur $S$, jointe au cas de caractéristique $> 0$
appliqué aux fibres, donne le résultat.\footnote{Une part importante des pages
73 et 74 est illisible et la transcription le dit. La structure de l'argument
--- étalement sur un schéma de type fini sur $\mathbb{Z}$, puis descente du cas
$p > 0$ aux fibres --- est ce que les phrases portantes autorisent à
restituer ; le détail du contrôle de la dimension ne l'est pas, et il n'est pas
inventé ici.}

\subsection*{Les deux corollaires et la question}

\textbf{Corollaire 1.} Invariance birationnelle de
$\underline{\mathrm{Pic}}^{\tau}_{X/k}$ pour $X$ propre et lisse sur $k$.

\textbf{Corollaire 2.} Soit $K$ une extension séparable primaire de type fini
de $k$ admettant un modèle propre et \emph{lisse} sur $k$. Alors
$\varinjlim_{X} \underline{\mathrm{Pic}}^{\tau}_{X/k}$, $X$ parcourant les
modèles propres et normaux de $K$ sur $k$, existe --- le système inductif étant
\emph{essentiellement constant}.

\textbf{Question.} Le corollaire 2 reste-t-il valable en supposant seulement
que $K$ admette un modèle \emph{normal} ? Cela équivaudrait à trois énoncés de
bornitude, sur tous les modèles propres normaux :
\begin{enumerate}
  \item $\dim \underline{\mathrm{Pic}}_{X/k}$ reste majoré ;
  \item l'ordre de la torsion de $\mathrm{Pic}(X/k)$ reste majoré ;
  \item $\dim H^{1}(X, \mathcal{O}_{X})$ reste majoré.
\end{enumerate}
Avec, en marge, l'observation que (i) résulte de (iii).

\subsection*{La borne uniforme, et la question finale}

\textbf{Théorème.} Soit $K$ une extension de type fini du corps $k$. Il existe
un entier $n(K/k)$ tel que, pour tout modèle $X$ de $K$ propre et normal sur
$k$,
\[
  \dim \underline{\mathrm{Pic}}_{X/k} \;\leqslant\; n(K/k) .
\]

\emph{Démonstration.} On se ramène au cas où $k$ est algébriquement clos, en
posant $\bar{K} = K \otimes_{k} \bar{k}$, extension de type fini de $\bar{k}$.
On utilise alors
\[
  \dim \underline{\mathrm{Pic}}_{X/k} \;=\; \tfrac{1}{2}\,
  \operatorname{rang}_{\mathbb{Z}_{\ell}}
  \bigl(\text{partie $\ell$-primaire de } \pi_{1}(X/k) \text{ rendu
  abélien}\bigr) ,
\]
et l'on conclut par le théorème de majoration uniforme du $\pi_{1}$ pour les
modèles normaux --- lequel résulte de la partie facile de Lefschetz, prise sur
les groupes fondamentaux rendus abéliens.\footnote{L'égalité est celle du
module de Tate : pour $X$ propre et normal sur $k$ algébriquement clos, la
partie libre du complété $\ell$-adique de $\pi_{1}(X)^{\mathrm{ab}}$ est
$T_{\ell}\bigl(\underline{\mathrm{Pic}}^{0}_{X/k}\bigr)$, de rang
$2\dim\underline{\mathrm{Pic}}^{0}$ ; et $\underline{\mathrm{Pic}}^{\tau}$ et
$\underline{\mathrm{Pic}}^{0}$ ont même dimension, leur quotient étant fini.
La page énonce l'égalité sans ces précisions ; c'est bien le rang de la partie
libre qui intervient, non celui du groupe entier.}

De façon plus précise : prenant un modèle $X$ projectif propre et normal, on
prend dans $X$ une « courbe générique $C$ » --- celle du lemme de Bertini
ci-dessus --- et pour $n(K)$ le \emph{genre} de cette courbe générique. C'est
donc, à un facteur près, la même borne que celle du groupe fondamental
birationnel : les deux moitiés de la chemise donnent le même entier.

\begin{tikzcd}
C' \arrow[r, hook] \arrow[d] & X' \arrow[d] \\
C_{k(u)} \arrow[r, hook] & X \arrow[r, hook] & \mathbf{P}^r
\end{tikzcd}

La chemise s'achève sur une question, et sur trois points d'interrogation en
marge : \emph{y a-t-il une démonstration du fait que cela donne une majoration,
n'utilisant pas la théorie du groupe fondamental et le fait que
$\underline{\mathrm{Pic}}^{\tau}_{X/k}$ est propre sur $k$ si $X$ est normale ?}

\pagerange{77}{80}
\section*{« Revêtements formels étales et théorème de Mumford » : Murre revient (pages 77 à 80)}

Un feuillet de sa main porte le titre, et sous lui vient une lettre de Murre
datée du 30 avril 1969 --- deux ans après l'échange sur la ramification
modérée. Elle accompagne la \emph{version préliminaire} des notes que Murre a
rédigées sur les revêtements formels modérément ramifiés et sur un résultat de
« plumbing », et elle demande un avis avant d'aller plus loin : « To me the
entire thing seems somewhat indigestible ». Le chapitre sur l'analogue du
théorème de Mumford reste à écrire.

Suivent deux pages de commentaires, chapitre par chapitre, dont trois méritent
d'être relevées parce qu'elles disent où la théorie fait mal.

\begin{itemize}
  \item Chapitre I : les morphismes étales de schémas formels sont définis
    comme systèmes inductifs de morphismes étales de schémas ordinaires, et le
    lemme crucial est que les voisinages étales formels sont localement des
    complétés de voisinages étales ordinaires.
  \item Chapitre III : « This is a nuisance ! » Murre n'a pas suivi le nouveau
    dispositif --- celui des données de ramification --- et il en donne la
    raison : « one gets the feeling that the notion of tamely ramified covering,
    although natural, is very unmanageable if one has to use \emph{the functor}
    instead of the geometric realization itself ». C'est la réponse de
    l'utilisateur au manifeste de la page 31, et elle mérite d'être lue à côté
    de lui : ce que Grothendieck tient pour indispensable à la compréhension,
    Murre le tient pour impraticable à la rédaction.
  \item Chapitre III encore : c'est en 3.6 que la réalisation géométrique
    \emph{détermine} le revêtement modérément ramifié, et Murre note que la
    normalité, plutôt que la régularité, suffit d'ordinaire --- avec cette
    difficulté que la normalité n'est peut-être pas stable par localisation
    étale formelle.
\end{itemize}

Le reste des commentaires est technique : quels lemmes servent où, quelle
hypothèse ($D_{i} \cap D_{j} = \varnothing$) est imposée par quel manque, et
l'aveu, pour le chapitre IX, que certains points ont été utilisés sans être
certains « because it seems to me that you used them implicitely and without
mentioning in your notes ». Les notes de lecture que Grothendieck a prises sur
ce texte sont aux pages 81, 83 et 85, dans le lot 5.

\end{document}
